% Generated draft: requires mathematical and source review.
\documentclass[11pt]{article}
\usepackage[margin=25mm]{geometry}
\usepackage{fontspec}
\setmainfont{Latin Modern Roman}
\usepackage{amsmath,amssymb,mathtools,unicode-math}
\setmathfont{Latin Modern Math}
\usepackage{enumitem,needspace,xurl,hyperref}
\setlength{\parindent}{0pt}
\setlength{\parskip}{5pt}
\setlength{\emergencystretch}{3em}
\newcommand{\N}{\mathbb{N}}
\newcommand{\Z}{\mathbb{Z}}
\newcommand{\Q}{\mathbb{Q}}
\newcommand{\R}{\mathbb{R}}
\newcommand{\C}{\mathbb{C}}
\newcommand{\F}{\mathbb{F}}
\newcommand{\A}{\mathbb{A}}
\newcommand{\PP}{\mathbb{P}}
\newcommand{\E}{\mathbb{E}}
\newcommand{\eps}{\varepsilon}
\newcommand{\dd}{\,\mathrm{d}}
\begin{document}
\section*{UnsolvedMath ID: 30006308}
\textbf{Source problem code:} catalogue entry
\subsection{Definitions and mathematical statement}
Let \(X\) be a smooth complete toric variety. The supplied record does not specify the ground field, its characteristic, the dimension of \(X\), or whether deformations are ordinary, equivariant, or subject to additional conditions. Completeness is part of the hypothesis; projectivity is not assumed.

For a possible formal interpretation, fix a field \(k\) and a deformation problem for \(X\) over \(k\). Suppose that a miniversal formal deformation base is described by a complete Noetherian local \(k\)-algebra \((R_X,\mathfrak m_X)\) with residue field \(k\). Its tangent cone at the distinguished point is then
\[
C_X=\operatorname{Spec}\!\left(\operatorname{gr}_{\mathfrak m_X}R_X\right),\qquad
\operatorname{gr}_{\mathfrak m_X}R_X
=\bigoplus_{n\geq 0}\mathfrak m_X^n/\mathfrak m_X^{n+1}.
\]
In this convention, the number of irreducible components of the base means the number of minimal prime ideals of \(R_X\). These are conditional conventions for interpreting the question, not specifications supplied by the source. For ordinary infinitesimal deformations of a smooth variety over a field, first-order deformation classes are described by
\[
T_X^1=\operatorname{Ext}^1_{\mathcal O_X}(\Omega^1_{X/k},\mathcal O_X)
\cong H^1(X,\mathcal T_{X/k}),\qquad
\mathcal T_{X/k}=\mathcal H om_{\mathcal O_X}(\Omega^1_{X/k},\mathcal O_X).
\]
This description does not specify the simplicial complexes intended by the source.

\begin{enumerate}
\item Is the deformation space of every smooth complete toric variety determined by its tangent cone? Under the formal convention above, one possible precise interpretation asks whether, for any two such varieties \(X\) and \(Y\) over the same field and with the same deformation convention,
\[
\operatorname{gr}_{\mathfrak m_X}R_X\cong
\operatorname{gr}_{\mathfrak m_Y}R_Y
\quad\Longrightarrow\quad
R_X\cong R_Y,
\]
where the first isomorphism is of graded \(k\)-algebras and the second is of complete local \(k\)-algebras. The source does not specify this interpretation or distinguish it from asking whether each deformation germ is isomorphic to its tangent-cone germ. These are different questions, and neither is selected here as the definitive intended statement.
\item Is there a relation between the number of irreducible components of the deformation space and the number of distinct simplicial complexes corresponding to first-order deformations of \(X\)? The supplied record defines neither the complexes, their correspondence with deformation classes, nor the equivalence relation meant by ``distinct.'' It also proposes no particular equality, bound, or other relation. Consequently, this part remains a request for a relation rather than a fully specified mathematical proposition.
\end{enumerate}

Both questions are retained as source questions. Their unresolved conventions require human mathematical review before a definitive formal statement can be fixed.
\subsection{Short English statement}
For a smooth complete toric variety, the source asks whether its deformation space can be recovered from its tangent cone, and whether its component count is related to a count of simplicial complexes describing first-order deformations. The deformation category, recovery criterion, complexes, and proposed numerical relation remain unspecified. The supplied open label is dated curation, not an independently verified conclusion.
\subsection{Sources}
\begin{itemize}
\item Local UnsolvedMath registry: \texttt{data/unsolvedmath/problems.json}, with \texttt{catalog/problems.json} as the supplementary registry specified by the project. Supplied permanent UnsolvedMath ID: \(30006308\); exact source problem code: \texttt{catalogue entry}. Title: \emph{Deformation Spaces of Smooth Complete Toric Varieties}.
\item Primary citation supplied in the record: \emph{Toric Geometry} (2025), Oberwolfach Reports, pp.~883--940, DOI \texttt{10.4171/owr/2025/19}. \url{https://doi.org/10.4171/owr/2025/19}. The record attributes the questions to workshop open-problem discussions collected in this report; no individual proposer or exact passage page is supplied.
\item Supplied literature source: Ilten and Robins, \emph{Locally Trivial Deformations of Toric Varieties}. \url{https://arxiv.org/abs/2409.02824}.
\item Supplied literature source: Ilten and Turo, \emph{Deformations of smooth complete toric varieties}. \url{https://msp.org/ant/2020/14-4/ant-v14-n4-p04-s.pdf}.
\item Literature assessment date: \texttt{2026-08-22}. This is supplied dated curation, not independently verified. The supplied assessment labels the problem open, says that the 2025 report poses both questions, and reports combinatorial descriptions of first-order and higher obstructions and examples without a general answer to either question. The record explicitly requests independent verification of time-sensitive status and cited claims.
\item The supplied statement status is \texttt{corrected\_verified}; its accompanying qualification says that only clear formula delimiters, notation, grammar, and typography were repaired, with mathematical content and scope preserved. This registry label is not independent mathematical verification of this draft.
\end{itemize}
\end{document}
